
\section{From the Agora}
\dexteremph{Flancian}\\

\subsection*{Introduction}

Thank you Ishmael, my friend, for creating this zine and for inviting me to participate. Thank you also, reader (potential friend!), for being here with us and reading these words.

My name is not important although you may know it freely; suffice it for now to say I am a Flancian. In these few pages I will try to show you what this means to me by telling you about my goals and values, my projects and hopes and assorted dreams. 

I will try to do so through a series of fragments as you might find them in the Agora of Flancia. Some of the fragments are in English, some are in my native tongue of Spanish; I hope you can find some of them enjoyable or interesting, but feel free to skim or skip as needed (in what probably mounts to the eternal ideal contract between writers and readers). More crucially, be aware that I am limited in my language skills\footnote{I’ve been trying to learn German for years, with only limited success. But one of my favorite words
is Bruchstückhaft, meaning “having the property of Broken-piece-ness”. Like my writing.} and so, when you read an imperfect word or a clumsy turn of phrase, please substitute a better one where available.

What is this Agora, you may ask at this point; and what is Flancia? These are reasonable questions whose answers may become more or less evident with optional time and patience.

But first, here's some earnest poems in Spanish to illustrate and motivate this enterprise from the lyrical.

\fleurona

\subsection*{Gloria}

\widepoetry{

Gloria!\\

Gloria!\\

Gloria!\\

[[Gloria]] a las maravillas del universo!\\

(O Glory be to the light of consciousness...)\\
}

\fleuronb 

\subsection*{Fork}
\widepoetry{
Fork!\\

Fork!\\

[[Fork]]!\\

As we merge we'll fork!\\

(As we catch we'll throw...)\\

Like a Rainbow folding!\\

Like Maitreya flowing!\\
}

\fleurona

\subsection*{Agua}
\widepoetry{

Agua!

Agua!

Agua!

[[Agua]] para los sedientos\\

Pan, (y Flan...) para los hambrientos\\

(Sal para los hipotensos...)\\

Fuego para los que tienen frío\\

Techo para los que no tienen casa\\

Luz para los que leen\\

Ágora para los que escriben\\

Flancia para los oprimidos!\\
}

\fleuronb

\subsection*{Gone (tanka)}
\widepoetry{
Gone gone,

gone beyond\\

Everybody gone

to the other shore\\

Gone free,

gone kindly\\

If you have to go,

go kindly\\

Meditate!
}

\fleurona 

\subsection*{Las Jaras}
\widepoetry{
Ay! Quién fuera

Una flecha derecha

Volando hacia el corazón de Moloch\\

Ay! Quién fuera

Como una saeta

Volando hacia el corazón de Moloch

Llevo en mi carcaj:  tres jaras.\\

Las jaras, qué jaras?\\

Las de Avalokiteshvara,\\

Tara y Maitreya!\\
}

\fleuronb

\subsection*{Lady Burup}
\widepoetry{
Ay! Quién tuviera

Los ojos verdes de Lady Burup

Sus ojos negros

trasnochando\\

Sus bellas rayas!\\

Ay! Quién tuviera\\

Los ojos negros de Lady Burup\\

Sus ojos verdes 

tornasolados\\

Sus rayas bellas!\\
}

\fleurona 

\subsection*{Struck by lightning}
\widepoetry{
Sometimes I feel as if I’m struck by lightning;

Sometimes I catch fire.\\

I hope you know what I mean;

I can’t tell you what to feel.\\

But some of these days\\

I feel like slaying Moloch\\
}

\fleuronb

\subsection*{Flow}
\widepoetry{
As we fork

we'll merge\\

As we merge

we'll fork\\

As we merge 

we'll flow\\

Like a [[Rainbow folding]]\\

Like Maitreya glowing\\
}

\fleurona


\subsection*{An Agora in its intention}

If I had to sum up the Agora in its intention, if not yet in its current incarnation, I guess I could do worse than beginning as such: an Agora for the oppressed seems necessary, and of priority if there is to be only one, but I would rather we built many. Agoras of knowledge, maybe most of them; Agoras filled with kindness, hopefully all. For the greater good of humanity and friends; for our thriving together, and the reduction of suffering in the known universe. But also just for being; for playing and working and coordinating our projects.

The greater Agora, more than a particular instance, is a class of space; each instance can be digital or embodied, urban or rural. I see the Agora overall as the composition of all willing instances, as connected by bridges.

What does it take to build an Agora? I say any space that identifies as an Agora is one, as long as it's pro-social, well-meaning and kind-hearted. Like the Agora of Flancia.

For more on the Agora, see https://anagora.org.

\subsection*{About Flancia}

Flancia is an experiment in the space of protopias. It manifests itself as writing, code, music and images.

What is Flancia like? It is a place very much like our world, where people found ways to cooperate rationally on common goals for the benefit of beings. Where we built bridges; from the urban to the rural, from the analog to the digital, from my garden to yours; all with each other and more.

For more on Flancia, see https://flancia.org.

\fleuronb

\subsection*{Poesía generativa}
\widepoetry{
Maybe, just maybe

Everything will go\\

We will be free, free, free

free from suffering!

We will burn:

- burn, slowly

- burn, kindly\\

If you have to burn, burn kindly!\\

Meditate,\\

and go with love!\\
}
